\documentclass[10pt,a4paper]{amsart}
\usepackage[margin=19mm]{geometry}
\usepackage{amsmath,amssymb,mathtools}
\usepackage[scr=rsfs,cal=euler]{mathalfa}
\usepackage{xfrac}
\usepackage[unicode,hidelinks,bookmarks=false]{hyperref}
\newcommand{\ie}{i.e.\ }
\newcommand{\cf}{cf.\ }
\newcommand{\resp}{respectively\ }
\newcommand{\Subsection}{\subsection}
\providecommand{\nobreakdash}{\nobreak\hbox{-}}
\setlength{\emergencystretch}{2em}
\multlinegap0pt
\usepackage{siunitx}
\usepackage{siunitx}
\usepackage{siunitx}
\usepackage{siunitx}
\usepackage{xcolor}
\usepackage{siunitx}
\usepackage{tikz}
\usepackage{amssymb}
\usepackage{algorithm}\usepackage{algpseudocode}
\usepackage{float}
\newtheorem{lemma}{Lemma}
\title{Lotka-Sharpe Neural Operators \\ for Control of Population PDEs}
\author{Miroslav Krsti\'c, Iasson Karafyllis, Luke Bhan, Carina Veil}
\date{Source: arXiv:2604.03892v1 (CC BY 4.0)}
\begin{document}
\maketitle

\noindent\emph{Source and licence.} Lotka-Sharpe Neural Operators \\ for Control of Population PDEs. Version arXiv:2604.03892v1 (CC BY 4.0), distributed by arXiv under the Creative Commons Attribution 4.0 International licence, http://creativecommons.org/licenses/by/4.0/, as recorded in the arXiv licence metadata for this exact version and shown on the abstract page https://arxiv.org/abs/2604.03892v1. Full licence notice: \texttt{licenses/arxiv-2604.03892-CC-BY-4.0.txt}.\par\smallskip

\noindent\textit{Editorial scope.} Counted unit: the article's lemma lem-Delta\_u (source0.tex lines 2717--2724). The article's complete original proof of it is delivered with it (source0.tex lines 2727--2803, one \texttt{proof} environment of 77 lines). No mathematics is written, repaired or re-derived: the statement and the proof are the article's own lines, in the article's own notation, and no retained line is substituted or re-flowed.  Retained uncounted as the source-local context the proof needs: lines 2643--2662.  Nothing was omitted from any retained span, so the package carries no omission record.  No retained line contains a citation, so the wrapper carries no bibliography and no bibliography entry.  No retained line contains a figure environment, an image-inclusion command or drawing code, so no image is delivered and the package declares no asset dependency.  Editorial formatting only, in addition: an AMS \texttt{amsart} preamble that restates the article's own declarations and macros used by the retained lines, namely the environments \texttt{lemma} on separate counters, together with the standard packages those lines need.  The retained environments print in this document as Lemma 1 in order of appearance, whereas the article prints the same items with the numbering of its own full document.  The complete native source of the article as distributed in the arXiv e-print for this exact version is bundled unchanged as \texttt{sources/197843/source0.tex}.
\begin{lemma}[Lipschitz continuity of $\mathcal{G}_\gamma$, $\mathcal{G}_\kappa$, $\mathcal{G}_\pi$ in $\zeta$]
\label{lem-Lip-other}
Let $g, k \in C^0([0,A];\mathbb{R}_{\geq 0})$, $\mu \in C^0([0,A];\mathbb{R}_{\geq 0})$, and let $[\zeta_{\min},\zeta_{\max}] \subset \mathbb{R}_{\geq 0}$ be a compact interval. Then:

\smallskip
\noindent\textit{(i)} The map $\zeta \mapsto \mathcal{G}_\gamma(g,\zeta,\mu)$ is Lipschitz on $[\zeta_{\min},\zeta_{\max}]$ with constant
\begin{equation}
L_\gamma := A\|g\|_\infty e^{\zeta_{\max} A}.
\end{equation}

\noindent\textit{(ii)} The map $\zeta \mapsto \mathcal{G}_\kappa(k,\mu,\zeta)$ is Lipschitz on $[\zeta_{\min},\zeta_{\max}]$ with constant
\begin{equation}
L_\kappa := A^2\|k\|_\infty e^{\zeta_{\max} A}.
\end{equation}

\noindent\textit{(iii)} The map $\zeta \mapsto \mathcal{G}_\pi(k,\mu,\zeta)$ is Lipschitz on $[\zeta_{\min},\zeta_{\max}]$, in the sup norm, with constant
\begin{equation}
L_\pi := \frac{A^2}{2}\|k\|_\infty e^{\zeta_{\max} A}.
\end{equation}
\end{lemma}

\begin{lemma}[$\Delta_u$ bounded in terms of $|e_1|+|e_2|$]
\label{lem-Delta_u}
Fix $\delta>0$. For every $R\in(0,\min\{a,b\})$ there exists $C_{R}>0$ such that
\begin{equation}
|\Delta_u(\eta,e_1,e_2)|\le C_{R}(|e_1|+|e_2|)
\end{equation}
for all $\eta\in\mathcal T_{R}:=\{\eta:\ r(\eta)\le R\}$ and $ |e_1|+|e_2|\le\delta$.
\end{lemma}

\begin{proof}
Fix $\delta>0$ and $R\in(0,\min\{a,b\})$. Let $\eta\in\mathcal T_{R}=\{\eta:\,r(\eta)\le R\}$. Since $r(\eta)=\sqrt{\phi_1(\eta_1)^2+\phi_2(\eta_2)^2}$, with $\phi_1(\eta_1)=a(1-e^{-\eta_1})$ and $\phi_2(\eta_2)=b(e^{\eta_2}-1)$, and since $R<\min\{a,b\}$, it follows that
\begin{equation}
1-\frac{R}{a}\le e^{-\eta_1}\le 1+\frac{R}{a},\qquad
1-\frac{R}{b}\le e^{\eta_2}\le 1+\frac{R}{b}.
\end{equation}
Hence there exist constants $E_{1,R},E_{2,R}>0$ such that
\begin{equation}
e^{-\eta_1}\le E_{1,R},\qquad e^{\eta_2}\le E_{2,R},\qquad \forall\,\eta\in\mathcal T_{R}.
\end{equation}
Let $|e_1|+|e_2|\le\delta$. Since $\hat\zeta_i=\zeta_i-e_i$, the arguments $\zeta_1-e_1$ and $\zeta_2-e_2$ lie in compact intervals, and the mappings $\mathcal{G}_\gamma$, $\mathcal{G}_\kappa$, and $\mathcal{G}_\pi$ are Lipschitz in $\zeta$ on these intervals. Hence there exist constants $L_{\Gamma_i,\delta},L_{K_i,\delta},L_{P_i,\delta}>0$ such that
\begin{align}
|\Gamma_1(e_2)|\le&\; L_{\Gamma_1,\delta}|e_2|,\\
|\Gamma_2(e_1)|\le&\; L_{\Gamma_2,\delta}|e_1|,\\
|K_1(e_1)|\le&\; L_{K_1,\delta}|e_1|, \\ 
|K_2(e_2)|\le&\; L_{K_2,\delta}|e_2|, \\ 
|P_1(e_1)|\le&\; L_{P_1,\delta}|e_1|, \\ 
|P_2(e_2)|\le&\; L_{P_2,\delta}|e_2|.
\end{align}
Since the interval $[\zeta_1-\delta,\zeta_1+\delta]$ is compact and the map 
$\zeta \mapsto \mathcal{G}_\gamma(g_2,\zeta,\mu_1)$ is continuous and strictly positive, it attains a positive minimum on this interval. 
Hence there exists $\underline{\gamma}_2>0$ such that
\begin{align}
\label{eq-lower-gamma}
\gamma_2 + \Gamma_2(e_1) =\;& \nonumber  \mathcal{G}_\gamma(g_2,\zeta_1 - e_1,\mu_1) \\\ge\;&  \min_{\zeta \in [\zeta_1-\delta,\zeta_1+\delta]} \mathcal{G}_\gamma(g_2,\zeta,\mu_1) =: \underline{\gamma}_2 > 0.
\end{align}
Similarly, there exist $\underline{\kappa}_2>0$ and $\underline{\pi}_1>0$ such that
\begin{equation}
\label{eq-lower-bnd-kappa-pi}
\kappa_2+K_2(e_2)\ge \underline{\kappa}_2>0,\qquad
\langle \pi_{0,1},n_1\rangle+P_1(e_1)\ge \underline{\pi}_1>0.
\end{equation}
By Lemma \ref{lem-Lip-other} and the lower bound in \eqref{eq-lower-gamma}, the map $\zeta \mapsto \frac{1}{x_1^*(0)\mathcal{G}_\gamma(g_2,\zeta,\mu_1)}$ is Lipschitz on $[\zeta_1-\delta,\zeta_1+\delta]$, and hence there exists $L_{a,\delta}>0$ such that
\begin{equation}
|\hat a-a|\le L_{a,\delta}|e_1|.
\end{equation}
Using the lower bounds in \eqref{eq-lower-gamma}, \eqref{eq-lower-bnd-kappa-pi}and the Lipschitz continuity of $\mathcal{G}_\gamma$, $\mathcal{G}_\kappa$, and $\mathcal{G}_\pi$, it follows that there exist constants $L_{m_1,\delta},L_{m_2,\delta}>0$ such that
\begin{equation}
|\hat m_1-m_1|\le L_{m_1,\delta}|e_1|,\qquad
|\hat m_2-m_2|\le L_{m_2,\delta}|e_2|.
\end{equation}
By definition,
\begin{equation}
\Delta_{\mathrm{gain}}(\eta,e_1,e_2):= \hat S(\eta)-S(\eta),
%\Delta_{\mathrm{gain}}(\eta,e_1,e_2)=\hat S(\eta)-S(\eta),
\end{equation}
where
\begin{align}
\hat S(\eta)-S(\eta)=&\;(1+\varepsilon)(e_1-e_2)-\varepsilon(\hat a-a)\nonumber \\ &-(\hat m_1-m_1)e^{-\eta_1} +(1+\varepsilon)(\hat m_2-m_2)e^{\eta_2}.
\end{align}
Therefore, for all $\eta\in\mathcal T_{R}$,
\begin{align}
|\Delta_{\mathrm{gain}}(\eta,e_1,e_2)|
\le&\;
(1+\varepsilon)(|e_1|+|e_2|)+\varepsilon|\hat a-a|\nonumber \\ &+|\hat m_1-m_1|e^{-\eta_1}+(1+\varepsilon)|\hat m_2-m_2|e^{\eta_2} \nonumber\\
\le&
(1+\varepsilon)(|e_1|+|e_2|)+\varepsilon L_{a,\delta}|e_1|\nonumber \\ &+L_{m_1,\delta}E_{1,R}|e_1|+(1+\varepsilon)L_{m_2,\delta}E_{2,R}|e_2| \nonumber\\
\le&
L_{\mathrm{gain},R,\delta}(|e_1|+|e_2|),
\end{align}
for some constant $L_{\mathrm{gain},R,\delta}>0$. Finally, using
\begin{equation}
\Delta_u(\eta,e_1,e_2)=-e_2-(\hat a-a)+\beta\,\Delta_{\mathrm{gain}}(\eta,e_1,e_2),
\end{equation}
we obtain
\begin{eqnarray}
|\Delta_u(\eta,e_1,e_2)|
&\le&
|e_2|+|\hat a-a|+\beta|\Delta_{\mathrm{gain}}(\eta,e_1,e_2)| \nonumber\\
&\le&
|e_2|+L_{a,\delta}|e_1|+\beta L_{\mathrm{gain},R,\delta}(|e_1|+|e_2|) \nonumber\\
&\le&
C_{R,\delta}(|e_1|+|e_2|),
\end{eqnarray}
for all $\eta\in\mathcal T_R$, where $C_{R,\delta}:=\max\{1,L_{a,\delta}\}+\beta L_{\mathrm{gain},R,\delta}$. Since $\delta$ is fixed throughout the lemma, we write $C_R := C_{R,\delta}$, noting that $C_R$ depends on both $R$ and $\delta$.
% for all $\eta\in\mathcal T_{R}$, where $C_{R,\delta}:=\max\{1,L_{a,\delta}\}+\beta L_{\mathrm{gain},R,\delta}$. Since $\delta$ is fixed, we denote this constant by $C_{R}$. This proves the lemma.
\end{proof}

\end{document}
